%% ============================================================================
%% Section: Compilation -- The Causal Snapshot
%% Included from main.tex via %% ============================================================================
%% Section: Compilation -- The Causal Snapshot
%% Included from main.tex via %% ============================================================================
%% Section: Compilation -- The Causal Snapshot
%% Included from main.tex via %% ============================================================================
%% Section: Compilation -- The Causal Snapshot
%% Included from main.tex via \input{sections/compilation}
%% Faithful to GraphLoader.java, TopologicalSorter.java, CompiledCausalGraph.java,
%% RuntimeRegistry.java, DefaultCausalGraphRuntime.java.
%% ============================================================================
\section{Compilation: The Causal Snapshot}\label{sec:compile}

Structural queries are not answered against the stored property graph directly.
On \texttt{causal.compile} the graph is read once and \emph{compiled} into an
immutable in-memory snapshot that every subsequent query reads. Compilation is
the boundary between the mutable, transactional graph and the read-optimized
structure of \S\ref{sec:index}: it is where we pay the $O(V{+}E)$ scan once, so
that queries need not pay it again.

\paragraph{Loading}
The loader performs two passes inside a single read transaction. The first pass
materializes every \textsf{Variable} node with its \textsf{observed}/\textsf{latent}
flag and, when present, its linear-Gaussian parameters (intercept~$\mu_0$, noise
variance~$\sigma^2$); undeclared properties fall back to defaults ($\mu_0{=}0$,
$\sigma^2{=}1$). The second pass reads the \textsf{CAUSES} relationships with a
Cypher query that filters on \lstinline[style=causaldsl]|coalesce(r.active,true)|,
so only \emph{active} edges enter the snapshot and a curator can toggle an edge
in or out of the model without deleting it; the optional edge coefficient~$\beta$
(\textsf{weight}, default $1.0$) is read alongside for the quantitative layer.
Filtering in Cypher lets the DBMS use its own relationship access rather than
forcing a full relationship scan in Java.

\paragraph{Topological order and acyclicity}
The loader then computes a topological order of the variables with Kahn's
algorithm over the forward adjacency. This serves two purposes. It is the order
in which the index closures of \S\ref{sec:index} and the SEM recursions of
\S\ref{sec:quant} are filled, and it is the \emph{acyclicity check}: if the order
fails to cover every variable, the graph contains a directed cycle and is not a
DAG, and compilation aborts with an error rather than admitting an ill-formed
causal model. Enforcing the DAG invariant at compile time means no downstream
operator has to defend against cycles.

\paragraph{The immutable snapshot}
The result is a \emph{compiled causal graph}: variables, forward and reverse
adjacency, the topological order, and the SEM parameters, all deep-copied into
unmodifiable structures. Immutability is deliberate. Because the snapshot never
changes after construction, every reader sees a consistent view without locking,
and the accelerating index (\S\ref{sec:index}) can be built once against it and
cached rather than reconstructed per query. It is also the single source of
truth: derived facts (ancestors, adjustment sets) are exposed through
\texttt{causal.*} procedures over the snapshot rather than materialized back into
the graph as extra relationships, which would double the maintenance burden
under updates.

\paragraph{Runtime registry and consistency}
Neo4j instantiates a procedure class afresh for every call, so the snapshot
cannot live inside a procedure. It is held instead by a process-scoped registry
that owns the current snapshot in a single \texttt{volatile} reference.
Recompilation and invalidation are serialized, while all reads are lock-free:
a query captures the current reference once and works against that fixed
snapshot, so a concurrent recompile installs a new snapshot without ever exposing
a partially built one. This gives a simple, snapshot-isolated consistency model:
a query runs entirely against the model as of the compile that produced its
snapshot. When and how a new snapshot is installed---by full recompilation, or by
the incremental patch that avoids re-reading the whole graph---is the subject of
\S\ref{sec:incr}.

%% Faithful to GraphLoader.java, TopologicalSorter.java, CompiledCausalGraph.java,
%% RuntimeRegistry.java, DefaultCausalGraphRuntime.java.
%% ============================================================================
\section{Compilation: The Causal Snapshot}\label{sec:compile}

Structural queries are not answered against the stored property graph directly.
On \texttt{causal.compile} the graph is read once and \emph{compiled} into an
immutable in-memory snapshot that every subsequent query reads. Compilation is
the boundary between the mutable, transactional graph and the read-optimized
structure of \S\ref{sec:index}: it is where we pay the $O(V{+}E)$ scan once, so
that queries need not pay it again.

\paragraph{Loading}
The loader performs two passes inside a single read transaction. The first pass
materializes every \textsf{Variable} node with its \textsf{observed}/\textsf{latent}
flag and, when present, its linear-Gaussian parameters (intercept~$\mu_0$, noise
variance~$\sigma^2$); undeclared properties fall back to defaults ($\mu_0{=}0$,
$\sigma^2{=}1$). The second pass reads the \textsf{CAUSES} relationships with a
Cypher query that filters on \lstinline[style=causaldsl]|coalesce(r.active,true)|,
so only \emph{active} edges enter the snapshot and a curator can toggle an edge
in or out of the model without deleting it; the optional edge coefficient~$\beta$
(\textsf{weight}, default $1.0$) is read alongside for the quantitative layer.
Filtering in Cypher lets the DBMS use its own relationship access rather than
forcing a full relationship scan in Java.

\paragraph{Topological order and acyclicity}
The loader then computes a topological order of the variables with Kahn's
algorithm over the forward adjacency. This serves two purposes. It is the order
in which the index closures of \S\ref{sec:index} and the SEM recursions of
\S\ref{sec:quant} are filled, and it is the \emph{acyclicity check}: if the order
fails to cover every variable, the graph contains a directed cycle and is not a
DAG, and compilation aborts with an error rather than admitting an ill-formed
causal model. Enforcing the DAG invariant at compile time means no downstream
operator has to defend against cycles.

\paragraph{The immutable snapshot}
The result is a \emph{compiled causal graph}: variables, forward and reverse
adjacency, the topological order, and the SEM parameters, all deep-copied into
unmodifiable structures. Immutability is deliberate. Because the snapshot never
changes after construction, every reader sees a consistent view without locking,
and the accelerating index (\S\ref{sec:index}) can be built once against it and
cached rather than reconstructed per query. It is also the single source of
truth: derived facts (ancestors, adjustment sets) are exposed through
\texttt{causal.*} procedures over the snapshot rather than materialized back into
the graph as extra relationships, which would double the maintenance burden
under updates.

\paragraph{Runtime registry and consistency}
Neo4j instantiates a procedure class afresh for every call, so the snapshot
cannot live inside a procedure. It is held instead by a process-scoped registry
that owns the current snapshot in a single \texttt{volatile} reference.
Recompilation and invalidation are serialized, while all reads are lock-free:
a query captures the current reference once and works against that fixed
snapshot, so a concurrent recompile installs a new snapshot without ever exposing
a partially built one. This gives a simple, snapshot-isolated consistency model:
a query runs entirely against the model as of the compile that produced its
snapshot. When and how a new snapshot is installed---by full recompilation, or by
the incremental patch that avoids re-reading the whole graph---is the subject of
\S\ref{sec:incr}.

%% Faithful to GraphLoader.java, TopologicalSorter.java, CompiledCausalGraph.java,
%% RuntimeRegistry.java, DefaultCausalGraphRuntime.java.
%% ============================================================================
\section{Compilation: The Causal Snapshot}\label{sec:compile}

Structural queries are not answered against the stored property graph directly.
On \texttt{causal.compile} the graph is read once and \emph{compiled} into an
immutable in-memory snapshot that every subsequent query reads. Compilation is
the boundary between the mutable, transactional graph and the read-optimized
structure of \S\ref{sec:index}: it is where we pay the $O(V{+}E)$ scan once, so
that queries need not pay it again.

\paragraph{Loading}
The loader performs two passes inside a single read transaction. The first pass
materializes every \textsf{Variable} node with its \textsf{observed}/\textsf{latent}
flag and, when present, its linear-Gaussian parameters (intercept~$\mu_0$, noise
variance~$\sigma^2$); undeclared properties fall back to defaults ($\mu_0{=}0$,
$\sigma^2{=}1$). The second pass reads the \textsf{CAUSES} relationships with a
Cypher query that filters on \lstinline[style=causaldsl]|coalesce(r.active,true)|,
so only \emph{active} edges enter the snapshot and a curator can toggle an edge
in or out of the model without deleting it; the optional edge coefficient~$\beta$
(\textsf{weight}, default $1.0$) is read alongside for the quantitative layer.
Filtering in Cypher lets the DBMS use its own relationship access rather than
forcing a full relationship scan in Java.

\paragraph{Topological order and acyclicity}
The loader then computes a topological order of the variables with Kahn's
algorithm over the forward adjacency. This serves two purposes. It is the order
in which the index closures of \S\ref{sec:index} and the SEM recursions of
\S\ref{sec:quant} are filled, and it is the \emph{acyclicity check}: if the order
fails to cover every variable, the graph contains a directed cycle and is not a
DAG, and compilation aborts with an error rather than admitting an ill-formed
causal model. Enforcing the DAG invariant at compile time means no downstream
operator has to defend against cycles.

\paragraph{The immutable snapshot}
The result is a \emph{compiled causal graph}: variables, forward and reverse
adjacency, the topological order, and the SEM parameters, all deep-copied into
unmodifiable structures. Immutability is deliberate. Because the snapshot never
changes after construction, every reader sees a consistent view without locking,
and the accelerating index (\S\ref{sec:index}) can be built once against it and
cached rather than reconstructed per query. It is also the single source of
truth: derived facts (ancestors, adjustment sets) are exposed through
\texttt{causal.*} procedures over the snapshot rather than materialized back into
the graph as extra relationships, which would double the maintenance burden
under updates.

\paragraph{Runtime registry and consistency}
Neo4j instantiates a procedure class afresh for every call, so the snapshot
cannot live inside a procedure. It is held instead by a process-scoped registry
that owns the current snapshot in a single \texttt{volatile} reference.
Recompilation and invalidation are serialized, while all reads are lock-free:
a query captures the current reference once and works against that fixed
snapshot, so a concurrent recompile installs a new snapshot without ever exposing
a partially built one. This gives a simple, snapshot-isolated consistency model:
a query runs entirely against the model as of the compile that produced its
snapshot. When and how a new snapshot is installed---by full recompilation, or by
the incremental patch that avoids re-reading the whole graph---is the subject of
\S\ref{sec:incr}.

%% Faithful to GraphLoader.java, TopologicalSorter.java, CompiledCausalGraph.java,
%% RuntimeRegistry.java, DefaultCausalGraphRuntime.java.
%% ============================================================================
\section{Compilation: The Causal Snapshot}\label{sec:compile}

Structural queries are not answered against the stored property graph directly.
On \texttt{causal.compile} the graph is read once and \emph{compiled} into an
immutable in-memory snapshot that every subsequent query reads. Compilation is
the boundary between the mutable, transactional graph and the read-optimized
structure of \S\ref{sec:index}: it is where we pay the $O(V{+}E)$ scan once, so
that queries need not pay it again.

\paragraph{Loading}
The loader performs two passes inside a single read transaction. The first pass
materializes every \textsf{Variable} node with its \textsf{observed}/\textsf{latent}
flag and, when present, its linear-Gaussian parameters (intercept~$\mu_0$, noise
variance~$\sigma^2$); undeclared properties fall back to defaults ($\mu_0{=}0$,
$\sigma^2{=}1$). The second pass reads the \textsf{CAUSES} relationships with a
Cypher query that filters on \lstinline[style=causaldsl]|coalesce(r.active,true)|,
so only \emph{active} edges enter the snapshot and a curator can toggle an edge
in or out of the model without deleting it; the optional edge coefficient~$\beta$
(\textsf{weight}, default $1.0$) is read alongside for the quantitative layer.
Filtering in Cypher lets the DBMS use its own relationship access rather than
forcing a full relationship scan in Java.

\paragraph{Topological order and acyclicity}
The loader then computes a topological order of the variables with Kahn's
algorithm over the forward adjacency. This serves two purposes. It is the order
in which the index closures of \S\ref{sec:index} and the SEM recursions of
\S\ref{sec:quant} are filled, and it is the \emph{acyclicity check}: if the order
fails to cover every variable, the graph contains a directed cycle and is not a
DAG, and compilation aborts with an error rather than admitting an ill-formed
causal model. Enforcing the DAG invariant at compile time means no downstream
operator has to defend against cycles.

\paragraph{The immutable snapshot}
The result is a \emph{compiled causal graph}: variables, forward and reverse
adjacency, the topological order, and the SEM parameters, all deep-copied into
unmodifiable structures. Immutability is deliberate. Because the snapshot never
changes after construction, every reader sees a consistent view without locking,
and the accelerating index (\S\ref{sec:index}) can be built once against it and
cached rather than reconstructed per query. It is also the single source of
truth: derived facts (ancestors, adjustment sets) are exposed through
\texttt{causal.*} procedures over the snapshot rather than materialized back into
the graph as extra relationships, which would double the maintenance burden
under updates.

\paragraph{Runtime registry and consistency}
Neo4j instantiates a procedure class afresh for every call, so the snapshot
cannot live inside a procedure. It is held instead by a process-scoped registry
that owns the current snapshot in a single \texttt{volatile} reference.
Recompilation and invalidation are serialized, while all reads are lock-free:
a query captures the current reference once and works against that fixed
snapshot, so a concurrent recompile installs a new snapshot without ever exposing
a partially built one. This gives a simple, snapshot-isolated consistency model:
a query runs entirely against the model as of the compile that produced its
snapshot. When and how a new snapshot is installed---by full recompilation, or by
the incremental patch that avoids re-reading the whole graph---is the subject of
\S\ref{sec:incr}.
